\section{Introduction}

\subsection{Resiliency Real-Time and Certification}

\begin{frame}{Introduction}
    Integrating AI models into systems that requiring predictable and safe behavior:
    \begin{itemize}
        \item \textbf{Challenge:} Black-box nature of AI models.
        \item \textbf{Issue:} Missing of structured standards and methodologies.
        \item \textbf{Starting point:} The \href{https://dl.acm.org/doi/full/10.1145/3626314}{Artificial Intelligence for Safety-Critical Systems in Industrial and Transportation Domains: A Survey} paper.
        \item \textbf{Repository:} \href{https://github.com/edoardosarri24/AI-Safety-Critical-runtime-monitoring}{GitHub}.
    \end{itemize}
\end{frame}

\begin{frame}{Introduction}
    \begin{block}{Terminology}
        \begin{itemize}
            \item \textbf{AI and Machine Learning:} Algorithmic paradigm that allows models to learn patterns from data.
            \item \textbf{AI Safety:} Methodologies to deploy AI in safety-critical systems in compliance with safety standards, avoiding/mitigating catastrophic harm.
            \item \textbf{Functional Safety (FuSa):} \textit{Part of the overall safety of a system that assures the freedom from unacceptable risk} (ISO 26262). Explicitly, is the characteristics of the system that bring back and mantain the system in a safe state.
        \end{itemize}
    \end{block}
\end{frame}